


Numbers are in \texttt{results/results.json} (keys given) and in \texttt{docs/03\_results/REDESIGN\_A\_RESULTS.md}, \texttt{REDESIGN\_B\_RESULTS.md}, \texttt{XSIM\_VERDICTS.md}. Every table rests on the constants of \texttt{docs/03\_results/ASSUMPTIONS.md}. Where an answer is a limit, it is stated as one.

\textbf{1. Is the simulator valid?}
Not shown. The custom simulator is an abstraction (static shadowing, a capture rule, an assumed loss model). It was compared with ns-3 (lorawan module) on 16 baseline cells: 15 agree within 5 points, one is 5.4 points out, so the pre-stated criterion was \textbf{not met} (\texttt{XSIM\_CALIBRATION.md}, \texttt{rd.x1.*}). The calibration was also run with both simulators at ns-3's shipped (7 dB better) gateway sensitivity: the criterion is not met there either (largest difference 7.9 points, \texttt{rd.x1.b.*}). On the D1 comparison the two simulators give different verdicts for H-D1a and H-D1b (\texttt{XSIM\_VERDICTS.md}, \texttt{rd.x13.agree.*}); D1's gain is therefore reported as simulator-dependent, and by the rule fixed in advance D1 does not go to the main paper. No result is validated against a real radio. Open: the comparison with the three published simulator papers (X1.4) was not run.

\textbf{2. Is D1 (or D2) just prior work?}
Possibly in part. Retransmission limits and ACK timeouts are in the LoRaWAN rules and in Khan et al. (full text read); an adaptive-retransmission paper for mines (IEEE CNS 2020) could not be opened and may overlap with D1; aggregated and group ACKs listing devices exist for LoRaWAN (A2S2-LoRaWAN, DG-LoRa; abstract level) so D2 does not own the aggregation. What D1 adds is stated as a combination (per-class deadlines, duty budget, Beta loss estimate, load index, one offline table), not as a unique idea (\texttt{docs/06\_references/CLOSEST\_PRIOR\_WORK.md}, \texttt{NOVELTY\_SCAN.md}). Absence statements read "no study found in the sources checked" with the sources listed (IEEE Xplore, Scopus, Google Scholar and TRID were not run; OpenAlex returned HTTP 503).

\textbf{3. There is no hardware.}
Correct. Every protocol number is SIMULATED or ANALYTICAL. The firmware decision path was run in a host build against the Python reference (0 mismatches in 10,000 random histories and in replayed simulator decisions, \texttt{rd.d6.*}); ESP32 flash, RAM, timing and an emulator run are NOT\_RUN. A bench protocol with logging hooks and a parser for the two-stage alert is ready (\texttt{docs/08\_hardware/R3\_BENCH\_PROTOCOL.md}); the hardware row is NOT\_RUN until the author supplies logs.

\textbf{4. Are the baselines strong enough?}
Six baselines are used and every D-hypothesis against the best of them per cell where stated (H-D1s): the severity table, UA-SCT, fixed 3 retries, quick retries, Tamang-style slots, Khan-style retries. On FRESH-A with the six baselines D1 beat the best one by +0.031 on average but only narrowly in some cells (\texttt{rd.fa.d1s.*}). A Nessa-style baseline (adaptive ACK timeout from the airtime plus a redundant retransmission; implemented from a secondary description because the CNS 2020 paper could not be opened; tuned on the DEV grid only; ADDED-AFTER-SESSION-1) was then added: D1 is better than it by +0.050 on average and in 75\% of the cells, but it is the best of the seven baselines in 17 cells and H-D1s with seven baselines is not met (\texttt{rd.fc.*}). D1 is better than quick retries by about 3.6 points, and quick retries are what most of D1's gain over the weak severity table consists of. On the cross-simulator scenario H-D1s is not met in either simulator. Standard LoRaWAN confirmed uplink (NbTrans 8, ACK in RX1 or RX2) was added in ns-3 (\texttt{rd.x13.ns3.mean\_delivery.lorawan}); ADR on was not run; no adaptive scheme from the literature could be implemented from the text read. The severity table is a weak baseline in the custom simulator (it waits 5 s before the initial Critical send); this is stated in \texttt{REDESIGN\_A\_RESULTS.md}.

\textbf{5. Was anything tuned on test data?}
No. Each redesign has three data roles (diagnosis on old grids, development on a new grid with at most 5 logged iterations, one fresh grid drawn from a hash after the code was frozen; \texttt{PREREGISTRATION\_V12.md} sections 1, 5, 10.5). The freeze files and grids are \texttt{results/tables/rd\_frozen\_A.json}, \texttt{rd\_fresh\_A\_grid.json}, \texttt{rd\_frozen\_B.json}, \texttt{rd\_fresh\_B\_grid.json}, committed before each run; each run writes a marker and refuses a second run. The DEV log is \texttt{results/tables/rd\_dev\_log.csv}. Disclosed limits: the simulator files changed after the FRESH-A freeze (D1 and D1 + D2 reproduce the stored delays bit for bit, \texttt{test\_rd\_regression\_a.py}); the cross-simulator run re-used the SF9 cells of FRESH-A (a cross-check, not a second held-out test); DEV and FRESH share the finite pool of real class pairs (V12 addendum A 5f).

\textbf{6. The data are synthetic.}
The simulator's events are drawn from real out-of-fold class pairs and real early-exit decision times of the VZCrash test events (MEASURED\_SOFTWARE pools); arrivals, geometry, losses and bursts are simulated. No synthetic crash pulse is used: no generator passed the realism gate (\texttt{DIAGNOSIS\_SYNTH.md}), so D5 used no synthetic data. Crash rate (0.1 per vehicle-hour) and burst sizes are stress values, not real rates.

\textbf{7. Are the statistics sound?}
Paired bootstrap over events within cells (2,000 resamples, seed 42), made more conservative by also resampling simulated hours (the verdict uses the larger p and the smaller bound), an intersection-union rule for bounds that must hold in every cell, and one Holm family of ten hypotheses with final adjusted p-values (\texttt{REDESIGN\_B\_RESULTS.md}). Every effect has a 95\% interval. Limits: the hour-cluster bootstrap still treats hours as independent; one-sided criteria with margins were chosen before the runs; with 100 to 600 hours per cell the intervals are narrow and describe the simulator, not the world.

\textbf{8. Dataset leakage.}
Splits are by group (recording), with the dedup and strict-test checks of the earlier rounds (\texttt{docs/04\_audits/}, \texttt{docs/00\_overview/PROJECT\_STATUS.md}); VZCrash has no vehicle identifier, so the most independent grouping available is used and the choice is in \texttt{DECISIONS.md}. The redesign round uses stored test-split decisions and does not refit any model. The dataset events were pre-selected on the device; this is a stated limit.

\textbf{9. Can it be reproduced?}
Seed 42 and fixed seeds per grid; pinned requirements; the redesign commands are in \texttt{docs/05\_data\_code/REPRODUCE.md} and the module docstrings; the freeze hashes are in the freeze files. Environment limits: the ns-3 build is in WSL outside the repository (versions, commits and local patches in \texttt{XSIM\_BUILD.md}); the system Python (numpy 2.x) was used for the redesign work, with three old tests failing on a numpy attribute; ns-3 runs are not bit-identical across machines. About 0.2\% of ns-3 runs crashed inside the module and are dropped for all policies (\texttt{rd\_x13\_ns3\_dropped.json}).

\textbf{10. What are the deployment limits?}
No road test, no real radios, no ESP32 timing. Duty-cycle limits are parameters (the 10\% for 433 MHz in India is the limit of rule 3 row 1 of G.S.R. 347(E), 2022, read in full on 2026-11-16). Hostile conditions are checked only in bounded models (H-D2c: two vehicles, small bounds; a stolen gateway key is outside the claim). The D1 table is about 187 kB per base SF; the C path uses double precision and needs a float or fixed-point port for the ESP32. Vehicle-to-vehicle overhearing is rare in the simulated channel, which is why D3 had nothing to act on. A crash alert system built on these results would still need field tests.


\subsection*{Independent review of the final manuscript (final verification run, 2026-11-15)}

A reviewer who had not seen the analysis read \texttt{paper/main.tex} and the supplement as an IEEE Access reviewer and listed ten weak points. Document fixes were applied; the rest is answered here with a pointer to the evidence. The review is model output, so each point was checked against the files before action.

{\scriptsize\raggedright
\begin{longtable}{p{0.237\linewidth}p{0.237\linewidth}p{0.237\linewidth}p{0.237\linewidth}}
\toprule
\# & Point & Action & Answer and evidence \\
\midrule
\endhead
1 & The study is called end-to-end, but the chain is never run as one pipeline & Fixed & The title no longer says end-to-end; the introduction says that the components are evaluated separately on different data and that no pulse passes through the whole chain \\
2 & Several results are fixed by design (the 1.87 s earlier arrival, the retraction share, the simulated-log accuracies, the attack suite) & Fixed & Text now says so: earlier arrival and retraction share are design and detector properties repeated in every column; the simulated-log results are sanity checks of the code against the generator's assumptions; the attack suite is a regression check by the author on the author's code \\
3 & The retry conclusion is stronger than the evidence & Fixed & The finding is limited to single-link replays without contention; the ns-3 ranking of quick retries is stated next to it (\texttt{docs/03\_results/XSIM\_VERDICTS.md}, \texttt{results.json} keys \texttt{rd.x13.*}, \texttt{r3.t1.*}) \\
4 & The trace replay mapping from counter gaps to retry times is not justified & Fixed (text) & The text now states the mapping (attempt j sees the j-th next frame of the same device, PREREGISTRATION\_V14, replay rules), says it says nothing at a shorter time scale, and calls the sequences gateway-reception sequences of weak links. The mapping itself is a limit of the method and stays \\
5 & Preregistration is not independently time-stamped; hypotheses were added after results; no correction across versions; events may be dependent & Fixed (text), open (analysis) & Text says author-registered and committed in git, lists the hypotheses added after results, notes the dependence of events and that no correction across versions was made. A family-wise count across versions was not computed (it would be a new analysis); the ledger shows per-version families. The abstract no longer quotes the supported share \\
6 & The C1 claim is broader than the evidence & Fixed (text) & Text defines a window duplicate, says only test events were removed and the forest and thresholds were not refitted, and limits the result to this forest. Retraining on de-duplicated training data was not run (new experiment); \texttt{docs/04\_audits/VZCRASH\_AUDIT.md} lists the sets. The release and revision that credit the finding must be cited by the author (\texttt{paper/PLACEHOLDERS.md}) \\
7 & The C2 threshold (80\% within 10\%) is arbitrary; counts are mixed; one sentence on passed modes was unclear & Partly fixed & The passed-modes sentence is explained (the other mode passed the gate, but no real range exists to compare). The 80\% rule is preregistered (V5 and V13); a sensitivity analysis over 75\% and 85\% was not run (new analysis, listed as future work); the shares for every setting are in the heat map and \texttt{sevstudy.grid.*} \\
8 & The formal contribution is narrower than presented & Fixed (text) & Text states 'in the symbolic model', names the out-of-scope attacks (jamming, key provisioning, master-key compromise, base-reboot replay) and gives the flash-write baseline (persisting every alert counter in the car and the base). The shared-key design M0 is the repository's own earlier design (\texttt{docs/01\_protocol/PACKET\_FORMAT.md}), not a published one; the paper says so by calling it the shared cell key design \\
9 & Detection lacks baselines and uncertainty & Fixed (text), open (analysis) & Text says that no simple baseline was run for the compact-forest comparison and adds the positive predictive value consequence. Earlier baselines in the repository are in \texttt{docs/03\_results/CODE\_VERIFICATION.md} and the severity studies; no new baseline was run \\
10 & Related work, length and reporting & Partly fixed & The abstract is 247 words (it was 246 to 253 during editing; the reviewer's 300 does not match); the hardware statement is cut to the abstract, introduction, threats and conclusion; the 40\% bound is called a design budget; the extra airtime versus single-frame share is explained; the India band and the LoED region mismatch is stated. The search method stays in the related-work section (the supplement has the full log). Figure error bars for the trace panel were not added (the intervals are in \texttt{results.json}); a frame-layout table was not added (layout in \texttt{PACKET\_FORMAT.md} and \texttt{FORMAL\_FIELD\_CHECK.md}) \\
\bottomrule
\end{longtable}}


